\subsection{An actual primitive motive for two-row rectangular coefficients}
The coefficient identity used in this subsection is known:
if $N=nd$, $0\le r\le\lfloor N/2\rfloor$, and $p_r(n,d)$ counts
partitions of $r$ in an $n\times d$ rectangle, then
\begin{equation}\label{eq:rectprimitiveknown}
 g_r:=g((N-r,r),(d^n),(d^n))=p_r(n,d)-p_{r-1}(n,d),
 \qquad p_{-1}=0.
\end{equation}
For $r\ge1$ the inspected primary locator is Pak--Panova,
\emph{Bounds on certain classes of Kronecker and $q$-binomial
coefficients}, author manuscript dated 13 October 2016,
Lemma~1.3, p.~2 \cite{PPBounds}; the source preserves the earlier
attributions. For $r=0$ the identity is the multiplicity one of the
trivial representation in a Specht tensor square. We apply this
known identity and classical Lefschetz theory; no historical
novelty is claimed for the underlying geometric interpretation.

\begin{finding}[\fid{F31}: a primitive Grassmannian coefficient object]\label{f:grassprimitive}
Let $n,d\ge1$, $N=nd$, $0\le r\le\lfloor N/2\rfloor$, and
$X=\Gr(n,n+d)$ over $\F_q$. Write $L=c_1(\mathcal O_X(1))$
for the Pl\"ucker class in $H^2(X_{\bar{\F}_q},\Q_\ell(1))$,
where $\ell\nmid q$. The actual Weil-module cokernel
\begin{equation}\label{eq:rectprimitivecoker}
 C_r=\operatorname{coker}\bigl(
 L:H^{2r-2}(X_{\bar{\F}_q},\Q_\ell)(-1)
       \longrightarrow H^{2r}(X_{\bar{\F}_q},\Q_\ell)\bigr)
 \simeq\Q_\ell(-r)^{\oplus g_r}
\end{equation}
is pure of weight $2r$, with scalar geometric Frobenius $q^r$;
the source is zero for $r=0$.

For $V_j=H^{2j}(X_{\bar{\F}_q},\Q_\ell)(j)$ and $r\ge1$, put
\[
 A_r=L^{N-2r+1}:V_r\longrightarrow V_{N-r+1},\qquad
 B_r=L^{N-2r+2}:V_{r-1}\longrightarrow V_{N-r+1}.
\]
Then $B_r$ is invertible and the Frobenius-equivariant projector
\begin{equation}\label{eq:rectprimitiveprojector}
 E_r=\id_{V_r}-L B_r^{-1}A_r
\end{equation}
has image $\ker A_r$ and kernel $L V_{r-1}$. For $r=0$ set
$E_0=\id$. Untwisting identifies its image with $C_r$.
It is the realization of an actual rational Chow-motive summand
$\mathsf P_r\simeq(\mathbb L^r)^{\oplus g_r}$, where
$\mathbb L$ denotes the Lefschetz motive. Its matrix in the
Schubert basis need not be nonnegative or integral.
\end{finding}

\begin{proof}
We first specify the arithmetic Tate decomposition. Let $\mathbb G_m$
act on the ambient space with weights $1,2,\ldots,n+d$.
The fixed $n$-planes are the rational coordinate planes indexed by
$I=\{i_1<\cdots<i_n\}$. The tangent weights are $j-i$ for
$i\in I$, $j\notin I$, so the positive-weight dimension is
\[
 a_I=\sum_{s=1}^n(d+s-i_s).
\]
The sequence $\lambda_s=d+s-i_s$ is a partition in the
$n\times d$ rectangle, and the inverse assignment is
$i_s=d+s-\lambda_s$. Thus $a_I=j$ occurs exactly $p_j(n,d)$
times. The motivic Bia\l ynicki--Birula decomposition in
Brosnan \cite[Theorems~3.3 and~3.5, author PDF pp.~5--6]{Brosnan}
applies to this smooth projective variety and split action.
Its fixed pieces are rational points, hence it gives, over $\F_q$,
\[
 h(X)\simeq\bigoplus_{j=0}^{N}(\mathbb L^j)^{\oplus p_j(n,d)},
 \qquad
 H^{2j}(X_{\bar{\F}_q},\Q_\ell)
    \simeq\Q_\ell(-j)^{\oplus p_j(n,d)},\quad H^{2j+1}=0.
\]
In particular Frobenius on $V_j$ is the identity. This actual
Tate decomposition supplies arithmetic splitting; purity alone
would not exclude a unipotent arithmetic extension.

The Pl\"ucker bundle is ample and defined over $\F_q$, and
$X$ is smooth projective of dimension $N$. BBD Theorem~5.4.10,
p.~144, applies to $X\to\Spec\F_q$ and the pure perverse
sheaf $\Q_\ell[N]$. Hard Lefschetz gives the stated isomorphism
$B_r$. Since $A_rL=B_r$, the map $L:V_{r-1}\to V_r$ is injective.
Set $T_r=L B_r^{-1}A_r$. Then
\[
 T_r^2=L B_r^{-1}(A_rL)B_r^{-1}A_r=T_r,\qquad
 A_r E_r=0,\quad E_r L=0,\quad
 E_r|_{\ker A_r}=\id.
\]
Also $E_rv=0$ implies $v=T_rv\in L V_{r-1}$.
Consequently
\[
 V_r=L V_{r-1}\oplus\ker A_r.
\]
The quotient map restricts to an isomorphism from the primitive
kernel. Its dimension is $p_r-p_{r-1}=g_r$ by the cited identity.
All maps are Frobenius equivariant, including $B_r^{-1}$, so
untwisting gives \eqref{eq:rectprimitivecoker}. More explicitly
the primitive subspace in the original degree is
\[
 \ker\bigl(L^{N-2r+1}:H^{2r}(X_{\bar{\F}_q},\Q_\ell)
 \longrightarrow
 H^{2N-2r+2}(X_{\bar{\F}_q},\Q_\ell)(N-2r+1)\bigr).
\]
The quotient-to-kernel splitting is canonical once $L$ is fixed.
The twist $(-1)$ in \eqref{eq:rectprimitivecoker} is necessary:
the rational divisor class is invariant in $H^2(1)$.

For the motivic assertion use rational coefficients in the displayed
Tate decomposition. Morphisms between sums of a fixed Tate motive
are rational matrices, and degree-zero morphisms between distinct
Tate powers vanish. After the displayed degree pieces are normalized
by their Tate powers, multiplication by $L$ and its powers are such
matrices, with rational entries given by algebraic correspondences.
Hard Lefschetz makes $B_r$'s determinant nonzero, so its inverse
is a rational matrix and an actual morphism in this Tate category.
Thus $E_r$ is an algebraic idempotent; the idempotent completion
gives its motive image with the stated rank. Brosnan's \S2,
author PDF pp.~3--4, specifies the correspondence and projector
category used here. This proves algebraicity in the cellular Tate
case; it does not assume the general standard conjecture on
Lefschetz inverses.

Finally take $X=\Gr(2,4)$ and $r=2$. In its integral Schubert
cycle lattice the add-one-box Pieri rule gives
\[
 L\sigma_1=\sigma_2+\sigma_{11},\qquad
 L\sigma_2=L\sigma_{11}=\sigma_{21}.
\]
These are the ordinary Schubert specializations of the Pieri rule
\cite[Theorem~7.10, p.~16]{Pieri}. Here $A_2=(1,1)$, $B_2=(2)$,
and
\[
 E_2=\frac12
 \begin{pmatrix}1&-1\\-1&1\end{pmatrix}.
\]
Its image is the primitive Tate line. The lattices
$\Z(\sigma_2+\sigma_{11})$ and
$\Z(\sigma_2-\sigma_{11})$ have sum of index two in
$\Z\sigma_2\oplus\Z\sigma_{11}$. Thus the prescribed
primitive-versus-Lefschetz decomposition does not split this
integral lattice, and its rational projector has negative entries.
This does not prohibit every integral projector onto the primitive
line: $\bigl(\begin{smallmatrix}1&0\\-1&0\end{smallmatrix}\bigr)$
is an integral idempotent with the same image and a different kernel.
\end{proof}

The two rectangular factors in \fid{F31} have arbitrary height and
width. Its restriction is the two-row output $(nd-r,r)$: it covers
single-part padded tails, not IP's general multipart tails. Within
IP's padded-tail range $|\bar\lambda|\le md$, it applies when
$r\le\min(md,\lfloor nd/2\rfloor)$; its own construction requires
no IP size threshold. The Satake coefficient in \fid{F20} for the
same two rectangles, using $\GL_n\times\GL_n$, has weight
$2nd(n-1)$, whereas \fid{F31} has weight $2r$. A Tate twist can
align weights, but their equal ranks do not choose an isomorphism,
an operation-compatible basis, or the nonstandard path comparison.
This finding uses classical hard Lefschetz and a known coefficient
identity; it does not derive that identity or purity anew from the
sign-only path argument.

% Separate paragraph for Part III (move there when integrating if desired).
\paragraph{\fid{F31} in Grothendieck's language.}
The integral Grassmannian is the projective Quot scheme of
rank-$d$ locally free quotients of $\mathcal O_{\Z}^{n+d}$,
with its universal sequence and Pl\"ucker line. FGA, Expos\'e~221,
\S3, Theorems~3.1--3.2, pp.~259--260, supplies the representability
framework. Its split $\GL_{n+d}$ action and the cocharacter used
above belong to SGA~3, I, 4.7.2--3. The cokernel is an actual
Weil object and the projector an actual rational correspondence,
with effective class $g_r[\mathbb L^r]$. The inverse exists here
because of the full matrix Hom spaces of the Tate motive, while
the integral Schubert cone need not be preserved. Cote~16,
scans~35, 37 and~39, gives the audited Lefschetz/primitive-projector
vocabulary; these leaves are not asserted to contain this
Kronecker identification or its proof.

% New bibliography keys requested by this fragment:
% PPBounds: I. Pak and G. Panova, Bounds on certain classes of Kronecker
% and q-binomial coefficients, author manuscript dated 13 October 2016,
% Lemma 1.3, p.2; https://www.math.ucla.edu/~pak/papers/QBC23.pdf.
% Brosnan: P. Brosnan, On motivic decompositions arising from the method
% of Bialynicki-Birula, Invent. Math. 161 (2005), 91--111; inspected author
% PDF https://math.umd.edu/~pbrosnan/Papers/bb.pdf, Section2 pp3--4,
% Theorems3.3 and3.5 pp5--6. These are author-PDF page locators.
